\section{Related Work}\label{sec:related}

\algname{} sits at the intersection of five research lines.  We
position it against each.

\paragraph{Probabilistic symbolic execution.}
Geldenhuys, Dwyer, and Visser~\cite{geldenhuys2012pse} introduced
probabilistic symbolic execution (PSE), combining symbolic execution
with model counting (LattE, Barvinok~\cite{koppe2008barvinok}) to
compute exact path probabilities for affine path conditions over
uniform inputs.  Filieri et al.~\cite{filieri2015tacas} extended PSE
to structured data; Borges et al.~\cite{borges2014compositional}
introduced compositional model counting.  Subsequent work, surveyed in
Pasareanu's synthesis lectures~\cite{pasareanu2022lecture}, has
broadened the class of treatable path conditions and input
distributions.

PSE is exact where it applies; \algname{} is statistical and
produces $(1-\delta)$-coverage intervals.  On programs where PSE
applies---bounded linear-integer-arithmetic paths and tractable input
distributions---PSE strictly dominates.  \algname{} fills the gap
when programs have unbounded loops, non-linear arithmetic, or
non-uniform input distributions.  When every region in \algname{}'s
frontier admits a closed-form mass and the closure rule succeeds on
every leaf, the two approaches agree exactly.

\paragraph{Statistical and probabilistic model checking.}
PRISM~\cite{kwiatkowska2011prism} and Storm~\cite{hensel2022storm}
verify quantitative temporal properties on Markov chains and MDPs.
Statistical variants---PLASMA-Lab~\cite{boyer2013plasma}, the Modest
Toolset~\cite{hartmanns2014modest}, and Storm's SMC
front-end~\cite{schoolderman2024storm}---discharge the same
specifications by sampling the underlying model.  These tools operate
on dedicated modelling languages (PRISM, JANI, Modest); \algname{}
operates directly on Python source, treating the program itself as the
model.  The trade-off is complementary: \algname{} cannot reason about
nondeterminism or scheduler choices, while PMC tools cannot reason
about source-level loops without a hand-built model.

\paragraph{Probabilistic assertions.}
Sampson et al.~\cite{sampson2014passert} introduced
\emph{passerts}---probabilistic assertions in C programs verified by
sampling with Hoeffding or Bayesian bounds.  \toolname{}'s
\code{failure\_probability} API targets the same query class.  The
difference is the bound and the sampling strategy: passerts use plain
Hoeffding sampling over the whole input space, whereas \algname{}
stratifies by symbolic path condition and uses a variance-adaptive
anytime-valid bound~\cite{waudbysmith2024betting}.  On programs whose
path structure concentrates probability mass, \algname{} reaches the
same certified half-width with substantially fewer samples; on
programs with flat, well-mixed distributions the advantage narrows.

\paragraph{Model counting.}
Approximate model counting~\cite{chakraborty2013approxmc,yang2023approxmc6}
provides PAC-$(\varepsilon,\delta)$ counts for propositional
satisfiability problems.  Recent work has added formal certificates
in Isabelle~\cite{kiesl2024certified}.  \algname{} does not currently
invoke a model counter---it uses importance sampling to estimate mass
on non-axis-aligned regions---but a natural extension would delegate
the region-mass subproblem to a certified counter, retaining
$(1-\delta)$-coverage with tighter constants.

\paragraph{Anytime-valid concentration inequalities.}
Classical Wilson and Hoeffding bounds are valid only at a
\emph{fixed} sample size; applying them under data-dependent stopping
requires a union bound over sample counts.  The mixture-martingale
literature, culminating in Howard et al.~\cite{howard2021time},
provides time-uniform Chernoff bounds via non-negative
supermartingales.  Waudby-Smith and
Ramdas~\cite{waudbysmith2024betting} sharpen this with a closed-form
predictable-plug-in empirical-Bernstein (PrPl-EB) confidence sequence
that is variance-adaptive and matches the asymptotic width of
Bennett's inequality without requiring the variance to be known.
\algname{} uses PrPl-EB per leaf with Bonferroni correction over the
(bounded) leaf count.

\paragraph{Property-based testing.}
QuickCheck~\cite{claessen2000quickcheck} and
Hypothesis~\cite{maciver2019hypothesis} search for a single input that
violates a property---a different goal from bounding the failure rate.
We provide a thin adapter (\code{dise.integrations.hypothesis}) that
converts Hypothesis strategies into (D1) distributions, so the same
property test can drive both edge-case search and operational
reliability estimation.

\paragraph{Summary.}
\algname{}'s technical contribution lies between the exact PSE line,
which it does not displace where PSE applies, and the sampling-only
literature, which it improves upon via symbolic stratification.  No
prior algorithm combines SMT-driven adaptive partitioning, closed-form
mass computation on axis-aligned (D1) regions, and an anytime-valid
empirical-Bernstein bound in a single framework.
